%------------------------------------------------------------------------------%

\usepackage{calculo_varias_variaveis-1}

\usepackage{tikz}
\usetikzlibrary{angles}

\title{Coordenadas Polares}

%------------------------------------------------------------------------------%
\begin{document}

\addtitle

%------------------------------------------------------------------------------%
\section{Sistema de Coordenadas Polares}

%------------------------------------------------------------------------------%
\begin{frame}{Coordenadas Polares}
  \centering
  
\begin{tikzpicture}

	% Coordinates
	\coordinate (A) at (6,0);
	\coordinate (B) at (0,0);
	\coordinate (C) at (2.5,2.5);
	\coordinate (B') at (2.5,3.5);
	\coordinate (A') at (1.8,3.2);

	% Axis
	\draw[thick,-latex] (-1ex,0) -- (6,0) node [below] {$x$};
	\draw[thick,-latex] (0,-1ex) -- (0,5) node [left]  {$y$};

	\draw[thick] (0,0) -- (2.5,2.5) node[pos=0.6, above left] {$r$};
  \node at (0.8,0.3) {$\theta$};

	% Help Lines
	\draw[dashed] (0,2.5) -- (2.5,2.5) -- (2.5,0);
	\draw[dashed] (3.53,0) arc (0:90:3.53);

	% Angle
	\pic[draw, ->, thick, angle eccentricity=1.7] {angle = A--B--C};

\end{tikzpicture}

\end{frame}

%------------------------------------------------------------------------------%
\begin{frame}{Coordenadas Polares}
  Definição
  \begin{itemize}[<+->]
    \item Origem ou polo
    \item Raio inicial
    \item \emph{$r$ distância orientada} até a origem, \(r\in\R\)
    \item \emph{$\theta$ ângulo orientado} a partir do raio inicial
          (em radianos) \(\theta\in\R\)
    \item Coordenadas polares $(r, \theta)$
  \end{itemize}
\end{frame}

%------------------------------------------------------------------------------%
\begin{frame}{Equações em Coordenadas Polares}
  \onslide<+->{\(r=a\)}
  \tabto{20mm}
  \onslide<+->{circulo com centro na origem e raio \(\abs{a}\)}
  \\[9mm]

  \onslide<+->{\(\theta=\theta_0\)}
  \tabto{20mm}
  \onslide<+->{reta passando pela origem com ângulo $\theta_0$ com o raio inicial}
\end{frame}

%------------------------------------------------------------------------------%
\begin{frame}{Relação Entre Coordenadas Polares e Cartesianas}
\begin{columns}
\column{0.2\linewidth}
\begin{align*}
  \onslide<2->{x          & = r\cos(\theta) \\}
  \onslide<3->{y          & = r\sin(\theta) \\[9mm]}
  \onslide<4->{r^2        & = x^2 + y^2     \\}
  \onslide<5->{\tan\theta & = \frac{x}{y}}
\end{align*}
\column{0.7\linewidth}
  \centering
  \onslide<1->{\includegraphics[width=\linewidth]{relacao_polares_catesianas}}
\end{columns}
\end{frame}

%------------------------------------------------------------------------------%
\section{Exemplos}

%------------------------------------------------------------------------------%
\begin{question}
  Encontre todas as coordenadas polares associadas ao ponto
  \(
    \left(r, \theta\right) =
    \left(2, \frac{\pi}{6}\right)
  \)
\end{question}

%------------------------------------------------------------------------------%
\begin{resolution}{Solução}
  O ponto
  \(
    \left(r, \theta\right) =
    \left(2, \frac{\pi}{6}\right)
  \)
  também pode ser escrito como
  \[
    \left(r, \theta\right) =
    \left(\hphantom{-} 2, \hphantom{-5}\,\frac{\pi}{6}+2n\pi\right)
    \qquad n=0, \pm 1, \pm 2, \dots
  \]
  ou
  \[
    \left(r, \theta\right) =
    \left(-2, -\frac{5\pi}{6}+2n\pi\right)
    \qquad n=0, \pm 1, \pm 2, \dots
  \]
\end{resolution}

%------------------------------------------------------------------------------%

%------------------------------------------------------------------------------%
\begin{question}
  Converta as seguintes coordenadas cartesianas em coordenadas polares
  $(r, \theta)$

  a) \quad \((x, y) = \left(3, 4\right)\)

  b) \quad \((x, y) = \left(-2, 2\sqrt{3}\right)\)

  c) \quad \((x, y) = \left(0, -5\right)\)

  Apresente as coordenadas com os valores ``mais simples''
\end{question}

%------------------------------------------------------------------------------%
\begin{resolution}{Solução Item a}
  \[
    r^2    \;=\; x^2 + y^2
    \pause \;=\; 3^2 + 4^2
    \pause \;=\; 9 + 16
    \pause \;=\; 25
  \]
  \[
    \pause r = \pm 5
  \]

  \[
    \pause \tan\theta \;=\; \frac{y}{x}
    \pause            \;=\; \frac{4}{3}
  \]
  \[
    \pause \theta \;=\; \arctan\left(\frac{4}{3}\right) \pm n\pi
    \qquad n\in \Z
  \]
  \[
    \pause
    \left(r, \theta\right) =
    \left(5, \arctan\left(\frac{4}{3}\right)\right)
  \]
\end{resolution}

%------------------------------------------------------------------------------%
\begin{resolution}{Solução Item b}
  \[
    r^2    \;=\; x^2 + y^2
    \pause \;=\; (-2)^2 + \left(2\sqrt{3}\right)^2
    \pause \;=\; 4 + 4\times 3
    \pause \;=\; 16
  \]
  \[
    \pause r = \pm 4
  \]

  \[
    \pause \tan\theta \;=\; \frac{y}{x}
    \pause            \;=\; \frac{2\sqrt{3}}{2}
    \pause            \;=\; \sqrt{3}
  \]
  \[
    \pause \theta \;=\; \arctan\left(\sqrt{3}\right) \pm n\pi
    \pause        \;=\; \frac{\pi}{3} \pm n\pi
    \qquad n\in \Z
  \]

  \pause
  \[
    \left(r, \theta\right) =
    \left(4, \dfrac{2\pi}{3}\right)
  \]
\end{resolution}

%------------------------------------------------------------------------------%
\begin{resolution}{Solução Item c}
  \[
    r^2    \;=\; x^2 + y^2
    \pause \;=\; 0^2 + 5^2
    \pause \;=\; 25
  \]
  \[
    \pause r = \pm 5
  \]

  \[
    \pause \tan\theta \;=\; \frac{y}{x}
    \pause            \;=\; \frac{-5}{0}
    \pause            \quad \nexists
  \]
  \[
    \pause \theta \;=\; \frac{\pi}{2} \pm n\pi
    \qquad n\in \Z
  \]

  \[
    \pause
    \left(r, \theta\right) =
    \left(5, \dfrac{3\pi}{2}\right)
  \]
\end{resolution}

%------------------------------------------------------------------------------%

%------------------------------------------------------------------------------%
\begin{question}
  Converta as coordenadas polares em coordenadas cartesianas $(x, y)$

  a) \quad \((r, \theta) = \left( 4, \frac{ \pi}{6} \right) \)

  b) \quad \((r, \theta) = \left( 3, \frac{3\pi}{4} \right) \)

  c) \quad \((r, \theta) = \left( 5, \pi            \right) \)
\end{question}

%------------------------------------------------------------------------------%
\begin{resolution}{Solução Item a}
  \[
    x      \;=\; r \cos(\theta)
    \pause \;=\; 4 \cos\left(\frac{\pi}{6}\right)
    \pause \;=\; 4 \frac{\sqrt{3}}{2}
    \pause \;=\; 2 \sqrt{3}
  \]

  \[
    \pause y \;=\; r \sin(\theta)
    \pause   \;=\; 4 \sin\left(\frac{\pi}{6}\right)
    \pause   \;=\; 4 \frac{1}{2}
    \pause   \;=\; 2
  \]

  \[
    \pause (x, y) = \left( 2 \sqrt{3}, 2 \right)
  \]
\end{resolution}

%------------------------------------------------------------------------------%
\begin{resolution}{Solução Item b}
  \[
    x      \;=\; r \cos(\theta)
    \pause \;=\; 3 \cos\left(\frac{3\pi}{4}\right)
    \pause \;=\; \frac{-3\sqrt{2}}{2}
  \]

  \[
    \pause y \;=\; r \sin(\theta)
    \pause   \;=\; 3 \sin\left(\frac{3\pi}{4}\right)
    \pause   \;=\; \frac{3\sqrt{2}}{2}
  \]

  \[
    \pause (x, y) = \left(-\dfrac{3\sqrt{2}}{2}, \dfrac{3\sqrt{2}}{2}\right)
  \]
\end{resolution}

%------------------------------------------------------------------------------%
\begin{resolution}{Solução Item c}
  \[
    x      \;=\; r \cos(\theta)
    \pause \;=\; 5 \cos\left(\pi\right)
    \pause \;=\; 5(-1)
    \pause \;=\; -5
  \]

  \[
    \pause y \;=\; r \sin(\theta)
    \pause   \;=\; 5 \sin\left(\pi\right)
    \pause   \;=\; 0
  \]

  \[
    \pause (x, y) = (-5, 0)
  \]
\end{resolution}

%------------------------------------------------------------------------------%
%------------------------------------------------------------------------------%
\begin{question}
  Encontre a equação polar para o círculo \(x^2+(y-3)^2=9\)
\end{question}

%------------------------------------------------------------------------------%
\begin{resolution}{Solução}
\onslide<+->{Escrevendo a equação em coordenadas polares}
\begin{align*}
  \onslide<+->{x^2 + (y-3)^2                 & = 9 \\}
  \onslide<+->{x^2 + y^2 - 6y + 3^2          & = 9 \\}
  \onslide<+->{x^2 + y^2 - 6y                & = 0 \\}
  \onslide<+->{r^2 - 6 r \sin(\theta)        & = 0 \\}
  \onslide<+->{r \big(r - 6\sin(\theta)\big) & = 0}
\end{align*}
\onslide<+->{Portanto}
\[
  \onslide<+->{r=0}
  \onslide<+->{
  \qquad\text{ou}\qquad
  r = 6\sin(\theta)}
\]
\end{resolution}

%------------------------------------------------------------------------------%

%------------------------------------------------------------------------------%
\begin{question}
  Escreva a equação
  \[
    r = \frac{4}{2\cos(\theta)-\sin(\theta)}
  \]
  em coordenadas cartesianas
  e identifique seu gráfico
\end{question}

%------------------------------------------------------------------------------%
\begin{resolution}{Solução}
  Sabemos que
  \(
    x = r\cos(\theta) \quad
    y = r\sin(\theta)
  \)

  \[
    \pause
    r = \frac{4}{2\cos(\theta)-\sin(\theta)}
  \]
  \[
    \pause
    2r\cos(\theta) - r\sin(\theta) = 4
  \]
  \[
    \pause
    2x - y = 4
  \]
  \[
    \pause
    y = 2x - 4
  \]

  \pause
  O gráfico é uma reta
\end{resolution}

%------------------------------------------------------------------------------%

%------------------------------------------------------------------------------%
\begin{question}
  Mostre que a equação polar
  \[
    r = \frac{4}{2 - \cos\theta}
  \]
  representa uma cônica
\end{question}

%------------------------------------------------------------------------------%
\begin{resolution}{Solução}
\begin{columns}[t]
\column{0.3\linewidth}
  \begin{align*}
    \onslide<+->{r                             & = \frac{4}{2 - \cos\theta} \\}
    \onslide<+->{\left(2 - \cos\theta\right)r  & = 4                        \\}
    \onslide<+->{\left(2 - \frac{x}{r}\right)r & = 4                        \\}
    \onslide<+->{2r - x                        & = 4                        \\}
    \onslide<+->{2r                            & = 4 + x                    \\}
  \end{align*}
\column{0.6\linewidth}
  \begin{align*}
    \onslide<+->{4r^2                          & = (4 + x)^2                \\}
    \onslide<+->{4\left(x^2+y^2\right)         & = 16 +8x + x^2             \\}
    \onslide<+->{4x^2 +4y^2 -8x -x^2           & = 16                       \\[7mm]}
    \onslide<+->{3x^2 +4y^2 -8x                & = 16}
  \end{align*}
\end{columns}
\end{resolution}
%------------------------------------------------------------------------------%

%------------------------------------------------------------------------------%
\begin{question}
  Escreva a equação
  \[
    r = 2\cos(\theta) - \sin(\theta)
  \]
  em coordenadas cartesianas e simplifique
\end{question}

%------------------------------------------------------------------------------%
\begin{resolution}{Solução}
  \onslide<+->{Sabemos que
  \[
    \cos(\theta) = \frac{x}{r}
    \qquad\qquad
    \sin(\theta) = \frac{y}{r}
    \qquad\qquad
    r^2 = x^2 + y^2
  \]}
\begin{columns}
\column{0.3\linewidth}
  \begin{align*}
    \onslide<+->{r & = 2\cos(\theta) - \sin(\theta) \\}
    \onslide<+->{r & = 2\frac{x}{r} - \frac{y}{r}   \\}
    \onslide<+->{r & = \frac{2x-y}{r}}
  \end{align*}
\column{0.5\linewidth}
  \begin{align*}
    \onslide<+->{r^2               & = 2x-y \\}
    \onslide<+->{x^2 + y^2         & = 2x-y \\}
    \onslide<+->{x^2 + y^2 - 2x+ y & = 0}
  \end{align*}
\end{columns}
\end{resolution}

%------------------------------------------------------------------------------%

%------------------------------------------------------------------------------%
\begin{question}
  Escreva a equação
  \[
    r = \frac{5}{\sin(\theta)-2\cos(\theta)}
  \]
  em coordenadas cartesianas e simplifique
\end{question}

%------------------------------------------------------------------------------%
\begin{resolution}{Solução}
  \onslide<+->{Sabemos que
  \[
    x = r\cos(\theta)
    \qquad\qquad
    y = r\sin(\theta)
    \qquad\qquad
    r^2 = x^2 + y^2
  \]}
  \vspace{-1.3\baselineskip}
  \begin{align*}
    \onslide<+->{r = \frac{5}{\sin(\theta)-2\cos(\theta)}   &     \\}
    \onslide<+->{r\left(\sin(\theta)-2\cos(\theta)\right) = & \;5 \\}
    \onslide<+->{r\sin(\theta)-2r\cos(\theta)             = & \;5 \\}
    \onslide<+->{y-2x = & \; 5   \\}
    \onslide<+->{y    = & \; 2x + 5}
  \end{align*}
\end{resolution}

%------------------------------------------------------------------------------%

%------------------------------------------------------------------------------%
\begin{question}
  Escreva a equação
  \[
    x^2 + xy + y^2 = 1
  \]
  em coordenadas polares e simplifique
\end{question}

%------------------------------------------------------------------------------%
\begin{resolution}{Solução}
  \onslide<+->{Sabemos que
  \[
    x = r\cos(\theta)
    \qquad\qquad
    y = r\sin(\theta)
    \qquad\qquad
    r^2 = x^2 + y^2
  \]}
  \vspace{-\baselineskip}
  \begin{align*}
    \onslide<+->{x^2 + xy + y^2                                         & = 1 \\}
    \onslide<+->{x^2 + y^2 + xy                                         & = 1 \\}
    \onslide<+->{r^2 + r\cos\theta r\sin\theta                          & = 1 \\}
    \onslide<+->{r^2 + r^2\cos\theta \sin\theta                         & = 1 \\}
    \onslide<+->{r^2\left(1 + \cos\theta \sin\theta \right)             & = 1 &}
    \onslide<+->{r & = \frac{\pm 1}{\sqrt{1 + \cos\theta \sin\theta}}}
  \end{align*}
\end{resolution}
%------------------------------------------------------------------------------%
\begin{resolution}{Solução Alternativa}
  \begin{align*}
    \onslide<+->{r^2\left(1 + \cos\theta \sin\theta \right)   & = 1 \\}
    \onslide<+->{r^2\left(1 + \frac{\sin(2\theta)}{2} \right) & = 1 \\}
    \onslide<+->{r^2\left(2 + \sin(2\theta) \right)           & = 2 \\}
    \onslide<+->{r & = \frac{\pm 2}{\sqrt{2 + \sin(2\theta)}}}
  \end{align*}
\end{resolution}

%------------------------------------------------------------------------------%


%------------------------------------------------------------------------------%
\section{Lista Mínima}

%------------------------------------------------------------------------------%
\begin{frame}{Lista Mínima}
  Cálculo Vol. 2 do Thomas 12\textsuperscript{a} ed. -- Seção 11.3
  \begin{enumerate}
    \item Estudar o texto da seção
    \item Resolver os exercícios: 1, 2, 6, 7, 11-15, 27-31, 53-57, 68
  \end{enumerate}

  \vfill
  Atenção: A prova é baseada no livro, não nas apresentações
\end{frame}

%------------------------------------------------------------------------------%
\end{document}
%------------------------------------------------------------------------------%
